\documentclass[11pt,a4paper]{article}
\usepackage[T1]{fontenc}
\usepackage{lmodern}
\usepackage[margin=24mm,headheight=14pt]{geometry}
\usepackage{amsmath,amssymb,amsthm,mathtools}
\usepackage{booktabs,longtable,array}
\usepackage{xcolor,microtype,fancyhdr}
\usepackage{enumitem}
\usepackage{listings}
\usepackage{hyperref}
\definecolor{ink}{HTML}{17364C}
\definecolor{muted}{HTML}{526170}
\hypersetup{colorlinks=true,linkcolor=ink,urlcolor=ink,citecolor=ink,
 pdftitle={Diagonal Power Certificates for Diophantine Consistency},
 pdfauthor={AI-assisted research prepared for the ProveIt project}}
\urlstyle{tt}
\setlength{\parskip}{4pt}
\setlength{\parindent}{0pt}
\setlist{itemsep=2pt,topsep=4pt}
\pagestyle{fancy}
\fancyhf{}
\fancyhead[L]{\small\color{muted}Diagonal power certificates}
\fancyhead[R]{\small\color{muted}ProveIt research contribution}
\fancyfoot[C]{\small\thepage}
\renewcommand{\headrulewidth}{0.3pt}
\newtheorem{theorem}{Theorem}[section]
\newtheorem{lemma}[theorem]{Lemma}
\newtheorem{proposition}[theorem]{Proposition}
\theoremstyle{definition}
\newtheorem{definition}[theorem]{Definition}
\newcommand{\Z}{\mathbb Z}
\newcommand{\N}{\mathbb N}
\newcommand{\Con}{\operatorname{Con}}
\newcommand{\Crefut}{\operatorname{Ref}_{\mathrm{C222}}}
\newcommand{\ZFCrefut}{\operatorname{Ref}_{\mathrm{ZFC}}}
\newcommand{\R}{\mathcal R}
\newcommand{\Qdiag}{Q_{214}}
\newcommand{\file}[1]{\nolinkurl{#1}}
\newcommand{\bitlen}{\operatorname{bitlen}}
\newcommand{\height}{\operatorname{H}}
\lstset{basicstyle=\small\ttfamily,breaklines=true,columns=fullflexible,
 frame=single,framerule=0.3pt,rulecolor=\color{muted},xleftmargin=3pt,xrightmargin=3pt}

\begin{document}
\begin{center}
{\Large\bfseries\color{ink} Diagonal Power Certificates\par}
\vspace{3pt}
{\Large\bfseries\color{ink} for Diophantine Consistency\par}
\vspace{10pt}
{\large A 214-variable integer quartic with coefficient height 219,813\par}
\vspace{12pt}
{Research prepared for Vladimir Reshetnikov's ProveIt project\par}
{\small AI-assisted mathematical research and reproducible arithmetic audit\par}
\vspace{5pt}
{25 October 2026\par}
\end{center}

\begin{abstract}
We give an explicit polynomial $\Qdiag\in\Z[x_0,\ldots,x_{213}]$,
of degree four and coefficient height $219{,}813$, for which ZFC proves
that absence of an integer zero is equivalent to its usual arithmetized
consistency statement. The polynomial is the sum of squares of 153
displayed quadratic residuals; its collected expansion has 5,656 monomials.
The construction modifies ProveIt package 173's direct C222 consistency
encoding. Its growth certificate is replaced by
\[
 t=464+z^2,\qquad K=464t^2,\qquad F=K^K.
\]
The power is represented by the existing ordinary polynomial Pell
certificate: the already authenticated index $K$ also serves as the base.
A positive-coefficient estimate for the actual C222 formula codes proves
that this growth suffices, allowing one intermediate variable and one
quadratic residual to be removed. At 214 variables and degree four, the
new height is more than five million times smaller than that of package
173's existing 214-variable companion, at the cost of 170 additional
monomials. Earlier self-power constructions are explicitly credited.
This is a concrete tradeoff within a direct small-coefficient construction,
not a claim of global simplicity or a new universal-polynomial record.
The ordinary proof uses identified inherited Pell and logical calibration
theorems. Exact independent arithmetic checks accompany it; a new
end-to-end proof-assistant formalization is not claimed.
\end{abstract}

\section{The equation and the precise result}

\begin{definition}
For the 153 quadratic residuals printed in Appendix~\ref{app:residuals}, set
\begin{equation}\label{eq:main}
 \boxed{\quad \Qdiag(x_0,\ldots,x_{213})
       :=\sum_{j=0}^{152}\R_j(x_0,\ldots,x_{213})^2=0,
       \qquad (x_0,\ldots,x_{213})\in\Z^{214}.\quad}
\end{equation}
Every variable is independently unrestricted over the integers. No
positivity, divisibility, exponentiation, or omitted side condition is
part of the definition.
\end{definition}

The complete quadratic coefficient list is
\file{direct/diagonal214/equation.json}; the complete collected quartic is
\file{direct/diagonal214/expanded_polynomial.json}. The appendix is generated
from the former and checked against it. Thus \eqref{eq:main} specifies an
ordinary fixed polynomial, rather than a polynomial schema or an
exponential Diophantine equation. Powers with variable exponents below
occur only in its mathematical interpretation.

\begin{theorem}\label{thm:main}
Using the identified C222/Pell and C222/ZF/ZFC calibration theorems of
ProveIt, the following equivalence is provable in ZFC:
\begin{equation}\label{eq:target}
 \bigl(\neg\exists\mathbf x\in\Z^{214}\;\Qdiag(\mathbf x)=0\bigr)
 \quad\longleftrightarrow\quad \Con(\mathrm{ZFC}).
\end{equation}
The polynomial has exact degree four, 5,656 nonzero collected monomials,
and coefficient height $\height(\Qdiag)=219{,}813$.
\end{theorem}

Here $\Con(\mathrm{ZFC})$ means that no finite derivation of falsity exists
in the fixed effective first-order ZFC presentation used by ProveIt.
This is the syntactic consistency predicate, not an assertion of external
truth or existence of a transitive set model. The ordinary proof establishes
the positive existence equivalence
\begin{equation}\label{eq:positive}
 \exists\mathbf x\in\Z^{214}\;\Qdiag(\mathbf x)=0
 \quad\longleftrightarrow\quad \ZFCrefut;
\end{equation}
negation gives \eqref{eq:target}. No consistency assumption is required
to prove these equivalences. In particular, this report does not prove
that the displayed equation lacks solutions.

The meaning of ``using'' in the theorem is ordinary mathematical reliance
on stated earlier theorems, not an extra consistency hypothesis. The new
arithmetic changes are proved here. The inherited unbounded logical
calibration is credited and its formalization boundary is stated in
Section~\ref{sec:boundary}.

\subsection{What improves, and what does not}

The parent is the principal equation of package 173,
\file{173-crt-index-anchor-215}, at repository commit
\texttt{02a78a08603cbae248c19c23a75907c6c6bd6443}~\cite{parent173}.
All equations in the following comparison have degree four and use
unrestricted integer variables.

\begin{center}
\small
\begin{tabular}{lrrrr}
\toprule
Equation & Variables & Residuals & Monomials & Coefficient height\\
\midrule
Package 173, principal $P_{215}$ & 215 & 154 & 5,660 & 91,679\\
Package 173, companion $P_{214}$ & 214 & 153 & 5,486 & 1,099,511,627,852\\
\textbf{This report, $\Qdiag$} & \textbf{214} & \textbf{153} & \textbf{5,656} & \textbf{219,813}\\
\bottomrule
\end{tabular}
\end{center}

Relative to $P_{215}$, one variable and one residual disappear, while
height increases. Relative to the existing $P_{214}$, height falls by a
factor greater than five million, while support increases by 170
monomials, about $3.10\%$. Package 173 also has a 213-variable companion,
whose height is $2{,}417{,}851{,}639{,}242{,}452{,}488{,}945{,}685$.
The broader repository contains much smaller variable counts with other
degree and coefficient tradeoffs. Accordingly, ``simplest'' has no
established single optimum here. The defensible claim is the explicit
measured improvement in the table.

The general diagonal-growth idea $F=t^t$ already occurs in package 170
and its sources~\cite{prior170}. The new contribution is the source-specific
label majorant, the amplified diagonal $F=K^K$ with $K=464t^2$, its
application to the range-free C-anchor interface, and the fully emitted
low-height endpoint. Neither the general self-power idea nor the
inherited C-anchor construction is claimed as new.

\section{A literal five-row change}

The parent coordinate chart has been optimized by translations, so
historical variable names are not by themselves semantic definitions.
The coordinates that matter here are exactly
\begin{align*}
 t&=x_0, & F&=x_2, & C&=2x_{16}, & K&=x_{93},\\
 W&=x_{131}, & A&=x_{132}, & E&=x_{139}, & q_{\mathrm{pow}}&=x_{140},\\
 D_K&=x_{181}, & Y_K&=x_{182}, & X_K&=x_{183}+1, & z&=x_{185}.
\end{align*}
$A$ is the common signed Pell center, $C$ the actual encoded row-index
representative, and $W$ the shared frame. In particular, $x_{183}$ is
$X_K-1$, not $X_K$. The unchanged rows 126 and 127 say
\begin{align}
 1-AY_K-D_K+x_{183}&=0,\label{eq:paira}\\
 -1+2D_K-D_K^2+2D_Kx_{183}+Y_K^2&=0.\label{eq:pairb}
\end{align}
Substituting $X_K=x_{183}+1=D_K+AY_K$ in \eqref{eq:pairb} gives
$X_K^2-(A^2-1)Y_K^2=1$. The retained cap and orientation certificates
authenticate its canonical index as $K$.

The old semantic growth data were
\[
 t=128+4z^2,\qquad T_2=t^2,\qquad K=T_2^2=t^4,\qquad F=t^K.
\]
The new data are
\begin{equation}\label{eq:newgrowth}
 \boxed{t=464+z^2,\qquad K=464t^2,\qquad F=K^K.}
\end{equation}
To implement this change, retain coordinates $x_0,\ldots,x_{213}$,
delete $x_{214}=T_2$, and change only the following parent rows:
\begin{center}
\begin{tabular}{rl}
\toprule
Row & New quadratic residual\\
\midrule
96 & $1+K^2-2KA+E$\\
97 & $F-D_K-KY_K+E q_{\mathrm{pow}}$\\
128 & $\R^{\mathrm{old}}_{128}+t^2-K^2$\\
130 & $t-464-z^2$\\
152 & $K-464t^2$\\
\bottomrule
\end{tabular}
\end{center}
Delete old row 153. Old rows 152 and 153 were $x_{214}-t^2$ and
$K-x_{214}^2$. No other row used $x_{214}$; no other row uses the horizon
root $z$. All other retained residuals are identical to the parent.

Rows 96 and 97 are the same singleton Pell power certificate with base
changed from $t$ to $K$. The index is still $K$, so no new principal pair
is introduced. The actual square frame formerly contained the growth
base $t$ and the index $K$. It now contains $K$ in both roles. The change
$+t^2-K^2$ in row 128 performs exactly this replacement; deleting the
existing index occurrence would be incorrect. The independent reviewer
reconstructed all 35 frame atoms, including affine translations and the
two occurrences of $K$.

This transformation preserves existence of witnesses through their
finite-refutation meaning. It is not a pointwise substitution between
complete old and new zero sets: the horizons and huge growth values
change. Both directions must therefore be proved semantically.

\section{A certified bound for the actual formula labels}\label{sec:labels}

The source calculus is the same square-shell recoding of C222 used by
the parent. Its constructor polynomials are
\begin{align*}
 I(a,b)&=2\bigl((a+b)^2+a\bigr),& M(a,b)&=I(a,b)+1,\\
 A_{\mathrm{syn}}(a,b)&=2\bigl((a+b)^2+3a+4b+3\bigr),&
 E_{\mathrm{syn}}(a,b)&=A_{\mathrm{syn}}(a,b)+1.
\end{align*}
The syntax constructor $A_{\mathrm{syn}}$ is distinct from the Pell center
$A$. Let $J(a,b,c,d)$ be the sum of the nine unchanged axiom-code
expressions and the five expressions $a,b,I(a,b),A_{\mathrm{syn}}(a,b),1$.
Every possible nonzero local conclusion or premise label is bounded by
$J$. Write $j(s)=J(s,s,s,s)$. This is a mathematical majorant only: no
new equation variable or range request for $J$ is added.

\subsection{The finite arithmetic certificate}

The authenticated source \file{majorant/shell_source.py}~\cite{shell112} is 2,955 bytes,
with SHA-256 recorded in Appendix~\ref{app:reproduce}. Its actual
\file{codes_only} path is executed through a positive arithmetic DAG by
\file{majorant/check_label_majorant.py}. The DAG has 312 nodes. All
coefficients are nonnegative, and the exact degrees of its fourteen
summands, in source order, are
\[
 16,8,8,8,8,32,16,8,524288,1,1,2,2,0.
\]
Thus $\deg j=524288$. The decisive exact inequalities are
\begin{align}
 \bitlen\bigl(j(464)\bigr)&=5{,}692{,}913,\label{eq:bits1}\\
 \bitlen\bigl(464^{645888}\bigr)&=5{,}721{,}264.\label{eq:bits2}
\end{align}
In particular,
\begin{equation}\label{eq:onepoint}
 j(464)<2^{5692913}<2^{5721263}\le 464^{645888}.
\end{equation}

The finite certificate uses exact integer arithmetic with outward-rounded
dyadic bounds, not floating-point logarithms. For example, it supplies
\begin{align*}
 j(464)&\le 12187191352475505376\cdot2^{5692849},\\
 464^{645888}&\ge14263689019699129227\cdot2^{5721200}.
\end{align*}
The matching lower and upper bounds locate each integer in a unique binary
interval. The optional \file{--exact} check evaluates the full integers
and confirms the bit lengths; that check was also run for this package.
It is unnecessary to expand the degree-524288 polynomial.

\begin{lemma}[Positive-coefficient extrapolation]\label{lem:coeff}
Let $p(X)=\sum_{i=0}^{D}a_iX^i$ with $a_i\ge0$. If $t\ge c>0$,
$d\ge D$, and $p(c)<c^d$, then $p(t)<t^d$.
\end{lemma}
\begin{proof}
For each $i\le D$, $t^i=c^i(t/c)^i\le c^i(t/c)^D$. Hence
\[
 p(t)\le p(c)(t/c)^D<c^d(t/c)^D\le t^d.
\]
Clearing denominators gives an equivalent integer argument.
\end{proof}

Set
\begin{equation}\label{eq:ldef}
 c_*=464,\qquad d_*=3c_*^2=645888,\qquad L(t)=t^{d_*}.
\end{equation}
Since $d_*\ge524288$, Lemma~\ref{lem:coeff} and
\eqref{eq:onepoint} give $j(t)<L(t)$ for every integer $t\ge464$.
The retained independent parameter ranges give $0\le a,b,c,d<t$ at every
decoded local row. Monotonicity therefore yields
\begin{equation}\label{eq:alllabels}
 J(a,b,c,d)\le j(t-1)<j(t)<L(t).
\end{equation}
This is a uniform proof for all possible witnesses. The single finite
evaluation is a certificate for the hypothesis of an unbounded
positive-coefficient lemma, not evidence from sampling.

The same certificate finds that $c=462$ fails this particular sufficient
test, while $c=463$ succeeds. We choose the convenient even coefficient
464, which supplies the parity margin used below. This observation is
not a minimality proof for the chosen constant or for any Diophantine
construction.

\section{The amplified diagonal growth lemma}

\begin{lemma}\label{lem:power}
If $t\ge c\ge3$ are integers, then $c^t\ge t^c$.
\end{lemma}
\begin{proof}
At $t=c$ there is equality. For every $n\ge c$,
\[
 (1+1/n)^c\le(1+1/c)^c
 \le\sum_{r=0}^{c}\frac1{r!}<3\le c.
\]
The strict bound follows from $r!\ge2^{r-1}$ for $r\ge2$ and the finite
geometric sum. Thus $(n+1)^c<cn^c$. Multiplying the induction hypothesis
$c^n\ge n^c$ by $c$ proves $c^{n+1}\ge(n+1)^c$.
\end{proof}

\begin{proposition}\label{prop:growth}
The semantic data \eqref{eq:newgrowth} satisfy
\begin{equation}\label{eq:growthbound}
 F\ge L(t)^t=t^{645888t}.
\end{equation}
\end{proposition}
\begin{proof}
Put $c=c_*=464$. Since $t\ge c$ and $K=ct^2$,
\begin{align*}
 F=K^K
  &=(ct^2)^{ct^2}
    =(c^t)^{ct}\,t^{2ct^2}\\
  &\ge t^{c^2t+2ct^2}
   \ge t^{3c^2t}=L(t)^t.
\end{align*}
The first inequality uses Lemma~\ref{lem:power}; the second uses $t\ge c$.
\end{proof}

The useful design principle is that a very large authenticated index can
also be its own base. The parent already paid for the $K$-index Pell pair.
Increasing its numerical role, while bounding the actual source labels
more accurately, replaces two quadratic growth-index gates by one.
This changes the growth certificate without changing the logical
calculus or introducing a machine simulation.

\section{Why the growth bound is enough}\label{sec:separation}

The following sharp lemma is inherited from package 173. Its proof is
included because it explains exactly what the new certificate must supply.

\begin{lemma}[Finite-product separation]\label{lem:separation}
Let $t,L,F$ be positive integers, $L\ge t$ and $F\ge L^t$.
Let $U,V$ be finite multisets in $\{1,\ldots,L\}$, with $|U|\le t$
and $|V|\le2t+1$. If
\[
 \prod_{u\in U}(1+uF)=\prod_{v\in V}(1+vF),
\]
then $U=V$ with multiplicities.
\end{lemma}
\begin{proof}
First recover the factor counts. If they differ, set
$r=\min(|U|,|V|)\le t$. The larger product is at least $(F+1)^{r+1}$;
the smaller is at most $(1+LF)^r\le L^r(F+1)^r$. Equality would give
$F+1\le L^r\le F$, a contradiction. Thus $|U|=|V|=m\le t$.
The case $m=0$ is immediate.

Replace $F$ by an indeterminate $Z$. Every coefficient of either product,
in each degree $0\le k\le m$, is a positive integer at most
\[
 \binom{m}{k}L^k\le m^{m-k}L^k\le L^m\le F.
\]
Here $\binom{m}{k}=\binom{m}{m-k}\le m^{m-k}$. Differences of corresponding
coefficients have absolute value strictly less than $F$, including when
one original coefficient equals $F$.

If the difference polynomial is nonzero, let $r$ be its least nonzero
degree, with coefficient $b_r$. Equality at $Z=F$, divided by $F^r$ and
reduced modulo $F$, forces $F\mid b_r$, contradicting $0<|b_r|<F$.
The polynomials are identical. Unique factorization in $\mathbb Q[Z]$,
or the distinct rational roots $-1/j$, recovers every label multiplicity.
\end{proof}

The decoded local rules have conclusion $v_i$ and two premises $p_i,q_i$;
missing premises are padded by truth, whose code is zero. Put $w_j=1+Fj$,
so $w_0=1$. The retained pending-state update is
\[
 S_{i+1}w_{v_i}=S_iw_{p_i}w_{q_i}.
\]
The beta endpoints give the initial state $F+1$ and final state $1$.
The possible boundary case $y=F$ would make the initial residue zero;
the positive weights would propagate zero to the final state, a
contradiction. Hence $y>F$, all states are positive, and multiplication
and cancellation of the $t$ updates gives
\begin{equation}\label{eq:balance}
 \prod_{i<t}w_{v_i}=(1+F)\prod_{i<t}w_{p_i}w_{q_i}.
\end{equation}
After discarding truth factors, the left side has at most $t$ factors
and the right side at most $2t+1$. Equations \eqref{eq:alllabels} and
\eqref{eq:growthbound} allow Lemma~\ref{lem:separation} to be applied.
Consequently conclusion multiplicity minus premise multiplicity is
one at falsity label 1 and zero at every other positive label.

Take the closure of truth under the finitely many instantiated rules.
If falsity were outside it, each rule with an outside conclusion would
have an outside premise. Its contribution to the total outside-label
balance would be nonpositive; the total could not be one. This contradicts
the recovered multiplicities. The argument extracts a finite C222
refutation even if the selected rule multiset contains cycles or repeated
premises. No topological ordering is assumed.

\section{Soundness: every integer zero gives a refutation}\label{sec:soundness}

Assume $\Qdiag(\mathbf x)=0$ for an arbitrary integer tuple. Since it is
a sum of integer squares, each residual vanishes. Row 130 gives $t\ge464$,
and row 152 gives $K=464t^2>0$, with $K$ even. The required order of the
remaining deductions is important.

\subsection{Recover the canonical powers before postponed signs}

The unchanged frame rows give a large nondegenerate signed center $A$;
the frame contains every actual requested base, index, and output.
The parent's C-anchor theorem~\cite{parent173} classifies the anchor
pair and bounds all principal indices before the sign of the actual CRT
representative $C$ is known. Its squared cap and bounded-error orientation
then authenticate every principal pair, including the retained $K$ pair.
This theorem remains applicable: the modified frame has the same form
and still contains all actual atoms. In particular, both the growth base
and growth index are the already positive $K$.

Rows 96 and 97 are exactly the singleton power certificate for this
base and index. They therefore force $F=K^K$. Thus $F>0$ is even before
the original sign bootstrap is used. Since $C=2x_{16}$ is even and the
retained divisibility is $B\mid C+F+1$, $B$ is nonzero and odd. The
self-sized radix multiplier remains 2 modulo 4, so its first nontrivial
power has a base of absolute value at least two. The unchanged singleton
and pooled power arguments can now recover the other powers and their
signs in the original order. None of this assumes the desired semantic
power equality before applying its certificate.

\subsection{The unchanged range and local-row interface}

All nine independent local-column range requests remain: four formula
parameters, two pending states, two beta quotients, and the encoded row
index. The only derived-label range previously removed by package 173
is still absent. Section~\ref{sec:labels} supplies its needed upper bound.
The local Horn polynomial, both beta endpoint congruences, the binomial
range certificates, and every-prime CRT compatibility are unchanged.

The inherited sign and beta bounds give $y\ge F$ and
$x\ge2+(t+1)y$, where $x$ denotes the beta code (not a free new coordinate).
For the local polynomial let $R_0=2+t+x+y+F$. The parent's uniform bound
is
\[
 0\le\mathrm{Local}<R_0^{2^{23}}.
\]
The new growth is stronger than $F\ge t^t$, and $t\ge464>16$; hence
$x>2^{32}$ and $R_0<2x$. If $2^k\le x<2^{k+1}$ with $k\ge32$, then
$2^{23}(k+2)<2^k\le x$, so $\mathrm{Local}<2^x$. Thus the unchanged
prime-power selector still upgrades local congruences to exact zero
rows. This is the same unbounded local estimate, not a check at particular
witness values.

For the binomial-selector radius, put
$\Lambda_t=\operatorname{lcm}(1,\ldots,t)$. The uniform estimate
\begin{equation}\label{eq:radius}
 \Lambda_t^2\le(t!)^2\le t^{2t}\le F
\end{equation}
follows from Proposition~\ref{prop:growth}. It avoids any parity
restriction on $t$. An earlier, sharper estimate proved only for even
horizons is not being used.

Consequently the decoded local rows satisfy the update and endpoints
used in Section~\ref{sec:separation}. Product separation and the
finite Horn-closure argument yield a C222 refutation. We have proved
\begin{equation}\label{eq:sound}
 \exists\mathbf x\in\Z^{214}\;\Qdiag(\mathbf x)=0
 \quad\Longrightarrow\quad\Crefut.
\end{equation}

\section{Completeness: every finite refutation gives a zero}

Start with a finite square-shell C222 refutation. Choose $z\in\N$ so
large that $t=464+z^2$ exceeds its local length and every independent
source parameter. These horizons are cofinal. Pad with legitimate truth
steps as in the parent, and set $K=464t^2$, $F=K^K$. All source labels
and all finite pending-multiset states are now fixed.

\subsection{The parity change is harmless}

Although $t$ need not be even, $464=16\cdot29$ gives
\begin{equation}\label{eq:valuations}
 v_2(K)=4+2v_2(t),\qquad
 v_2(F)=K v_2(K)\ge4K>2t.
\end{equation}
In particular $F$ contains more than the entire power of two dividing
$t!$. Let $O$ be the odd part of $t!$. Choose a sufficiently large
positive odd multiple $u$ of $O$ and set $y=Fu$. Then $t!\mid y$,
$y>F$, and $y$ exceeds all genuine pending states. It also has the
required literal shape $F(2w+1)$.

The beta moduli $1+(i+1)y$ are pairwise coprime. Indeed a common divisor
of the $i$th and $j$th moduli divides $(j-i)y$ and is coprime to $y$;
it therefore divides $j-i$, which divides $t!$ and hence $y$. The
divisor is one. Ordinary finite CRT supplies the beta code, chosen
cofinally large in its class.

Now set $B=\binom{F+t}{t}$ and $N=B^x$. For each $1\le i\le t$,
\eqref{eq:valuations} implies $v_2(F+i)=v_2(i)$. The binomial product
formula gives $v_2(B)=0$, so $B$, $N$, and $N^2$ are odd. The
parity-free radius estimate \eqref{eq:radius} also holds. Thus both
directions of the inherited generalized binomial CRT theorem apply:
every row has a selector, and every prime dividing $B$ has a compatible
row at its required full prime-power precision~\cite{binomial151}.

\subsection{Choose the reused anchor in a noncircular order}

The local-column CRT conditions determine nine residue classes modulo
$N^2$, one for each independent column. Choose the eight ordinary
representatives first. Then choose the actual row-index representative
$C$ in its prescribed class, positive and even, greater than all other
requested nominal indices, including $K$. This is always possible:
$N^2$ is odd, so the permitted even values form a cofinal progression
of step $2N^2$.

Only after this choice define the self-sized radix multiplier, its power
outputs, the actual square frame $W$, and its positive canonical Pell
center $A$. The modulus $N^2$ was fixed before choosing $C$ and does not
depend on this later frame. Increasing $C$ preserves all local CRT
conditions and strengthens the self-sized radix bounds. This is the
cofinal-C argument of the parent, with the new index $K$ included in the
finite list to be dominated.

The parent's cap and orientation completeness proof now supplies every
auxiliary Pell witness. The even index $C$ provides the same odd
auxiliary quotient; its canonical ordinate dominates the retained finite
sum of the other ordinates. Nonnegative slacks retain their ternary
representations. The new first power request uses canonical base $K$
and index $K$, and output $F=K^K$, so its divisor quotient is integral.
All other singleton, pooled, range, and local witnesses are unchanged
in form. Set the new index wire to $464t^2$ and omit the old private
$t^2$ wire. This constructs a zero of every new residual. Therefore
\begin{equation}\label{eq:complete}
 \Crefut\quad\Longrightarrow\quad
 \exists\mathbf x\in\Z^{214}\;\Qdiag(\mathbf x)=0.
\end{equation}

\section{Logical calibration and the proof boundary}\label{sec:boundary}

Equations \eqref{eq:sound} and \eqref{eq:complete} prove the arithmetic
existence equivalence with the unchanged square-shell C222 calculus.
The remaining logical chain is
\begin{align*}
 \text{square-shell C222 refutation}
 &\ \Longleftrightarrow\ \text{original C222 derivation of code 1}\\
 &\ \Longleftrightarrow\ \text{ZF refutation}\\
 &\ \Longleftrightarrow\ \text{ZFC refutation}.
\end{align*}
The first arrow is the parent's syntax recoding, a bijection fixing truth
0 and falsity 1 and transporting guarded rule instances. It is not a
claim that square-shell arithmetic constructors literally equal the
original numerical constructors.

The two named logical endpoints are
\begin{quote}\small\ttfamily
SetTheory.BusyBeaver.ZFC.KEndpoint.\allowbreak c222\_iff\_zf\_refutation\\
SetTheory.BusyBeaver.ZFC.ConstructibleChoice.\allowbreak zf\_refutation\_iff\_zfc\_refutation.
\end{quote}
They are pinned at commit
\texttt{953fb4d82d4229f4efc5f9db4cc2bf38ae329bb4}~\cite{kend,choice}.
Copies of their source files are included as dependency excerpts, with
byte and Git-blob hashes. Their transitive import closure is not bundled
or newly compiled. The endpoint source file also contains machine
theorems, but the named C222/refutation endpoint has no machine argument.
This construction introduces neither a K215 computation nor a
specialization of Jones' universal polynomial.

The new proof uses finite syntax, induction, ordinary integer inequalities,
polynomial identities, finite CRT, and the identified Pell theorems.
These are formalizable in ZFC and compose with the inherited calibration
to prove \eqref{eq:positive} and Theorem~\ref{thm:main}. All reasoning is
uniform over the internally coded integers and finite derivations of
ZFC. It is not restricted to externally standard integers or proofs.
No appeal to omega-consistency or external well-foundedness of a model
is made.

The largest logical dependency not independently rederived here is the
full C222/ZF calibration, including the seed-to-model lower direction.
Its internally well-founded model construction must not be replaced by
an assumption that consistency yields an externally transitive model.
On the arithmetic side, the delicate inherited dependencies are signed
Pell classification, squared divisibility, abscissa comparison, bounded
orientation, and pooled-power error estimates. The retained parent note
\file{inherited/parent173/ORDINARY_PROOF_DEPENDENCY_BOUNDARY.md} identifies
their precise sources.

The result's status is therefore an ordinary mathematical proof with
explicit inherited dependencies and exact finite arithmetic certificates.
It is not a new end-to-end Lean or Rocq proof object. The supplied
checks establish literal finite statements about the emitted equation
and source majorant; they do not on their own prove the unbounded
Pell, CRT, or logical interpretation.

\section{Verification, scope, and further research}

The producer authenticates the parent bytes before any transformation.
Two separately implemented finite checkers import none of that producer's
polynomial code. Between them they verify every changed row, every
unchanged row, the complete shared frame, index-pair roles, cap and
orientation terms, coordinate liveness, every coefficient of the
collected quartic, and the exact source-bound label calculation.
The root checker also checks all printed residuals, regeneration of the
TeX appendix, deterministic reproduction of all six equation artifacts,
and rejection of corrupted versions of the five risk-bearing rows.
The standard-library checks passed both normally and with Python
assertions disabled. Exact huge-integer evaluation was also performed.

The degree is at most four because every residual is quadratic. It is
exactly four because $x_{185}$ occurs only in
$\R_{130}=x_0-464-x_{185}^2$; the coefficient of $x_{185}^4$ in
$\Qdiag$ is one. All 214 variables occur in the collected polynomial.
The residual list contains 1,040 nonzero terms and has coefficient
height 464. Collecting all squares, including cancellations of mixed
terms, gives the stated 5,656 monomials and height 219,813.

The useful next research questions are narrower than an unsupported
claim of global simplicity. Can another already authenticated index
also absorb a power base without increasing the shared frame cost?
Can the source label majorant be made cheaper by changing the coding
while preserving a fully identified logical calibration? Can a
different Horn aggregation reduce the number of independent CRT
columns? Each would require a complete soundness and completeness
argument; finite numerical reductions alone do not establish it.

\begin{thebibliography}{9}
\raggedright
\bibitem{parent173}
V. Reshetnikov, \emph{ProveIt, package 173: CRT-index anchor, 215/214/213
integer quartics}. Pinned repository commit
\texttt{02a78a08603cbae248c19c23a75907c6c6bd6443}.
\href{https://github.com/VladimirReshetnikov/ProveIt/tree/02a78a08603cbae248c19c23a75907c6c6bd6443/Computability/DiophantineConsistency/con-zfc-diophantine-forms/173-crt-index-anchor-215}{Repository package and ordinary proof notes}.

\bibitem{prior170}
V. Reshetnikov, \emph{ProveIt, package 170: quotient carriers,
diagonal-pool source proof}. Same repository snapshot.
\href{https://github.com/VladimirReshetnikov/ProveIt/blob/02a78a08603cbae248c19c23a75907c6c6bd6443/Computability/DiophantineConsistency/con-zfc-diophantine-forms/170-quotient-carriers/direct/diagonal_pool/inputs/DIAGONAL_SOURCE_PROOF.md}{Earlier self-power construction}.

\bibitem{shell112}
V. Reshetnikov, \emph{ProveIt, package 112 intake: square-shell C222 source},
\file{ConZFC_Shell_Reciprocal_660/code/shell_source.py}, commit
\texttt{953fb4d82d4229f4efc5f9db4cc2bf38ae329bb4}.
\href{https://github.com/VladimirReshetnikov/ProveIt/blob/953fb4d82d4229f4efc5f9db4cc2bf38ae329bb4/Computability/DiophantineConsistency/con-zfc-diophantine-forms/112-powered-saturation-round14/intake/ConZFC_Shell_Reciprocal_660_2026-10-24/ConZFC_Shell_Reciprocal_660/code/shell_source.py}{Pinned source file}.

\bibitem{binomial151}
V. Reshetnikov, \emph{ProveIt, package 151: factorial-free binomial CRT},
\file{report/main.tex}, same pinned logical/source commit.
\href{https://github.com/VladimirReshetnikov/ProveIt/blob/953fb4d82d4229f4efc5f9db4cc2bf38ae329bb4/Computability/DiophantineConsistency/con-zfc-diophantine-forms/151-factorial-free-binomial-crt/report/main.tex}{Generalized selector theorem and every-prime completeness}.

\bibitem{kend}
V. Reshetnikov, \emph{ProveIt, KEndpoint.lean}, commit
\texttt{953fb4d82d4229f4efc5f9db4cc2bf38ae329bb4}.
\href{https://github.com/VladimirReshetnikov/ProveIt/blob/953fb4d82d4229f4efc5f9db4cc2bf38ae329bb4/Computability/BusyBeaver/Lean/BusyBeaver/ZFC/KEndpoint.lean}{C222/ZF refutation endpoint}.

\bibitem{choice}
V. Reshetnikov, \emph{ProveIt, ConstructibleChoice.lean}, same commit.
\href{https://github.com/VladimirReshetnikov/ProveIt/blob/953fb4d82d4229f4efc5f9db4cc2bf38ae329bb4/Computability/BusyBeaver/Lean/BusyBeaver/ZFC/ConstructibleChoice.lean}{ZF/ZFC refutation transport}.
\end{thebibliography}

\clearpage
\appendix
\section{Reproduction and artifact identities}\label{app:reproduce}

Run the following from the extracted package directory. The Python
programs require only the standard library.
\begin{lstlisting}
python3 -B verify.py --rebuild
python3 -B -O verify.py --rebuild --exact-majorant
python3 -B verification/review_diagonal.py
pdflatex -interaction=nonstopmode -halt-on-error report.tex
pdflatex -interaction=nonstopmode -halt-on-error report.tex
\end{lstlisting}
The first two commands check equation data and regenerate its printable
definition; they do not compile a proof-assistant development. The last
two compile this PDF using standard LaTeX packages. The evaluator
\file{evaluate.py} accepts an explicit JSON list of 214 integer
coordinates and evaluates either the quartic or all residuals, with
an optional modulus; it is not a solution search or a consistency test.

\begin{center}
\small
\begin{tabular}{p{0.39\linewidth}p{0.54\linewidth}}
\toprule
Artifact & Role\\
\midrule
\file{report.tex}, \file{residuals.tex} & Complete paper and all literal residuals\\
\file{direct/diagonal214/} & Exact equation, full expansion, metrics, variable map, source delta\\
\file{direct/build_diagonal.py} & Deterministic producer from authenticated parent\\
\file{direct/principal/} & Archived principal parent equation\\
\file{majorant/} & Authenticated source, positive-DAG checker, certificate, proof\\
\file{verify.py} & Independent artifact and delta verification\\
\file{verification/} & Separate review, checker, and execution receipts\\
\file{inherited/} & Credited proof dependencies and logical endpoint excerpts\\
\file{SOURCES.md} & Source paths, commits, and provenance\\
\bottomrule
\end{tabular}
\end{center}

The exact SHA-256 identities are listed below. Line breaks in the
display are typographical only; concatenate each pair of half-lines.
\begin{quote}\small\ttfamily
Parent equation:\\
8d919a1e5d6a802765b6d010c0140fb1d\\
83c3d36d211b209b7a8d2fb3f3fb6f6\\[4pt]
New equation:\\
7ffda8ff6f609918dc3d1824a28ff718\\
8ec8c3ea2be83abb1e4bc5d465fd7894\\[4pt]
Collected expansion:\\
171c6c26f97c0c0846f2aa4c6a8e134d\\
b6d4d2d459874b1e3b224bce67a8bcc4\\[4pt]
Square-shell source:\\
59c68a354fe5cdbc04bde33720080db2\\
436cb51ae25e3485ac5ae45ff4230199
\end{quote}
The equation digest authenticates the complete quadratic coefficient
list, rather than a shortened display. Verification receipts report
these same identities.

\clearpage
\section{Complete quadratic residuals}\label{app:residuals}

The following are all 153 residuals, indexed from 0 to 152. Every
symbol $x_i$ is one of the independently quantified integers
$x_0,\ldots,x_{213}$. There are no external constants to instantiate
and no omitted rows. Equation \eqref{eq:main} is the sum of their squares.
Rows that wrap continue on the next aligned line.

% Generated from the canonical quadratic residuals. All 153 are printed.
\begingroup
\fontsize{8.6}{10.4}\selectfont
\setlength{\jot}{3pt}
\allowdisplaybreaks[4]
\begin{align*}
\mathcal R_{0} &= x_{0} -x_{0}x_{4} -x_{4} +x_{7} \\
\mathcal R_{1} &= -x_{2} -x_{4}x_{5} +x_{6}x_{7} \\
\mathcal R_{2} &= -1 +x_{1} -x_{6}x_{7} \\
\mathcal R_{3} &= -2x_{2}x_{213} -x_{3}x_{18} +x_{4} +2x_{16} \\
\mathcal R_{4} &= -2 -x_{8} -x_{8}^{2} -2x_{8}x_{9} -2x_{9} -x_{9}^{2} +x_{19} \\
\mathcal R_{5} &= -2 -x_{8} -x_{8}^{2} -8x_{8}x_{19} -4x_{8}x_{20} -8x_{19} -16x_{19}^{2} -16x_{19}x_{20} -4x_{20} \\*
&\quad -4x_{20}^{2} +x_{21} \\
\mathcal R_{6} &= 1 +2x_{21} -4x_{21}^{2} +x_{22} \\
\mathcal R_{7} &= -2 -x_{9}^{2} -2x_{9}x_{10} -x_{10}^{2} +x_{23} \\
\mathcal R_{8} &= -1 +2x_{8} -x_{8}^{2} -2x_{8}x_{10} +2x_{10} -x_{10}^{2} +x_{24} \\
\mathcal R_{9} &= -7 +33x_{8} -35x_{8}^{2} -96x_{8}x_{10} -24x_{8}x_{19} -4x_{8}x_{23} -24x_{8}x_{24} +16x_{9} \\*
&\quad -2x_{9}^{2} -8x_{9}x_{23} +32x_{10} -64x_{10}^{2} -32x_{10}x_{19} -32x_{10}x_{24} +4x_{19} \\*
&\quad -4x_{19}^{2} -8x_{19}x_{24} +12x_{23} -4x_{23}^{2} +8x_{24} -4x_{24}^{2} +x_{25} \\
\mathcal R_{10} &= -16 +14x_{8} -2x_{8}^{2} -8x_{8}x_{9} -8x_{8}x_{23} +8x_{9} +8x_{23} -4x_{25}^{2} +x_{26} -2x_{43} \\
\mathcal R_{11} &= -6 +16x_{10} -16x_{10}^{2} -16x_{10}x_{23} -16x_{10}x_{24} +8x_{23} -4x_{23}^{2} -8x_{23}x_{24} \\*
&\quad +8x_{24} -4x_{24}^{2} +x_{27} \\
\mathcal R_{12} &= -5 +24x_{10} -16x_{10}^{2} -16x_{10}x_{19} -32x_{10}x_{24} -16x_{10}x_{27} +10x_{19} -4x_{19}^{2} \\*
&\quad -16x_{19}x_{24} -8x_{19}x_{27} +24x_{24} -16x_{24}^{2} -16x_{24}x_{27} +12x_{27} -4x_{27}^{2} \\*
&\quad +x_{28} \\
\mathcal R_{13} &= -12 -29x_{8} -37x_{8}^{2} -24x_{8}x_{9} -112x_{8}x_{10} -4x_{8}x_{23} -24x_{8}x_{24} -39x_{9} \\*
&\quad -37x_{9}^{2} -112x_{9}x_{10} -24x_{9}x_{23} -4x_{9}x_{24} -80x_{10} -128x_{10}^{2} -32x_{10}x_{23} \\*
&\quad -32x_{10}x_{24} -12x_{23} -4x_{23}^{2} -8x_{24} -4x_{24}^{2} +x_{29} \\
\mathcal R_{14} &= -18 -10x_{8} -2x_{8}^{2} -24x_{8}x_{9} -32x_{8}x_{10} -8x_{8}x_{23} -72x_{9} -72x_{9}^{2} \\*
&\quad -192x_{9}x_{10} -48x_{9}x_{23} -96x_{10} -128x_{10}^{2} -64x_{10}x_{23} -24x_{23} -8x_{23}^{2} \\*
&\quad -4x_{29}^{2} +x_{30} \\
\mathcal R_{15} &= -32 -42x_{8} +8x_{8}x_{10} +16x_{8}x_{23} -84x_{9} +16x_{9}x_{10} +32x_{9}x_{23} +21x_{10} \\*
&\quad -3x_{10}^{2} -4x_{10}x_{19} -8x_{10}x_{23} +21x_{19} -8x_{19}x_{23} +24x_{23} -4x_{23}^{2} +x_{31} \\*
&\quad -x_{59} \\
\mathcal R_{16} &= -1 -24x_{8} -64x_{8}^{2} -160x_{8}x_{10} -64x_{8}x_{23} -32x_{8}x_{24} -32x_{8}x_{31} -26x_{10} \\*
&\quad -100x_{10}^{2} -80x_{10}x_{23} -40x_{10}x_{24} -40x_{10}x_{31} -8x_{23} -16x_{23}^{2} \\*
&\quad -16x_{23}x_{24} -16x_{23}x_{31} -6x_{24} -4x_{24}^{2} -8x_{24}x_{31} -4x_{31} -4x_{31}^{2} +x_{32} \\
\mathcal R_{17} &= 1 -4x_{19}^{2} +x_{20} \\
\mathcal R_{18} &= -2 -x_{8} -x_{8}^{2} +8x_{8}x_{19} -4x_{8}x_{20} +8x_{19} -16x_{19}^{2} +16x_{19}x_{20} -4x_{20} \\*
&\quad -4x_{20}^{2} +x_{33} \\
\mathcal R_{19} &= 1 +2x_{33} -4x_{33}^{2} +x_{34} \\
\mathcal R_{20} &= -x_{9}^{2} -4x_{9}x_{34} -4x_{34}^{2} +x_{35} \\
\mathcal R_{21} &= -8 -20x_{8} -16x_{8}^{2} -64x_{8}x_{9} +32x_{8}x_{19} -40x_{9} -64x_{9}^{2} +64x_{9}x_{19} +18x_{19} \\*
&\quad -16x_{19}^{2} +x_{36} \\
\mathcal R_{22} &= 8x_{9} -62x_{9}^{2} -92x_{9}x_{10} -32x_{9}x_{23} -32x_{9}x_{36} +6x_{10} -34x_{10}^{2} \\*
&\quad -24x_{10}x_{23} -24x_{10}x_{36} -4x_{23}^{2} -8x_{23}x_{36} +4x_{36} -4x_{36}^{2} +x_{37} \\
\mathcal R_{23} &= 8x_{8} -64x_{8}^{2} -96x_{8}x_{10} -32x_{8}x_{24} -32x_{8}x_{37} +6x_{10} -36x_{10}^{2} \\*
&\quad -24x_{10}x_{24} -24x_{10}x_{37} +2x_{24} -4x_{24}^{2} -8x_{24}x_{37} +4x_{37} -4x_{37}^{2} +x_{38} \\
\mathcal R_{24} &= -x_{10}^{2} -2x_{10}x_{11} -x_{11}^{2} +x_{39} \\
\mathcal R_{25} &= -32 +22x_{8} -3x_{8}^{2} -8x_{8}x_{9} -8x_{8}x_{23} -8x_{8}x_{24} -4x_{8}x_{39} +16x_{9} -8x_{9}x_{24} \\*
&\quad -2x_{10}^{2} -8x_{10}x_{39} +16x_{23} -8x_{23}x_{24} +22x_{24} -4x_{24}^{2} +4x_{39} -4x_{39}^{2} \\*
&\quad +x_{40} -x_{43} \\
\mathcal R_{26} &= 6 -16x_{8}^{2} -36x_{8}x_{10} -32x_{8}x_{24} -40x_{8}x_{39} -16x_{8}x_{40} -22x_{10}^{2} \\*
&\quad -32x_{10}x_{24} -48x_{10}x_{39} -16x_{10}x_{40} -2x_{24} -16x_{24}^{2} -32x_{24}x_{39} -16x_{24}x_{40} \\*
&\quad +4x_{39} -24x_{39}^{2} -16x_{39}x_{40} -4x_{40}^{2} +x_{41} \\
\mathcal R_{27} &= x_{10} -25x_{10}^{2} -40x_{10}x_{24} -20x_{10}x_{27} -16x_{24}^{2} -16x_{24}x_{27} -4x_{27}^{2} \\*
&\quad +x_{42} \\
\mathcal R_{28} &= -1 +8x_{9} -4x_{9}^{2} -8x_{9}x_{23} +8x_{23} -4x_{23}^{2} +x_{43} \\
\mathcal R_{29} &= -1 -16x_{9}^{2} -8x_{9}x_{10} -32x_{9}x_{23} -16x_{9}x_{43} -x_{10}^{2} -8x_{10}x_{23} -4x_{10}x_{43} \\*
&\quad -16x_{23}^{2} -16x_{23}x_{43} -4x_{43}^{2} +x_{44} \\
\mathcal R_{30} &= 3 +x_{9} -2x_{9}x_{10} -4x_{9}x_{44} -3x_{10}^{2} -8x_{10}x_{44} -x_{23} -4x_{44}^{2} +x_{45} \\
\mathcal R_{31} &= -1 -4x_{45} -4x_{45}^{2} +x_{46} \\
\mathcal R_{32} &= -1 -x_{10}^{2} -8x_{10}x_{45} -4x_{10}x_{46} -16x_{45}^{2} -16x_{45}x_{46} -4x_{46}^{2} +x_{47} \\
\mathcal R_{33} &= 2 +x_{9} -x_{9}^{2} -4x_{9}x_{10} -4x_{9}x_{47} -4x_{10}^{2} -8x_{10}x_{47} -4x_{47}^{2} +x_{48} \\
\mathcal R_{34} &= -4x_{42}^{2} -8x_{42}x_{48} -4x_{48}^{2} +x_{49} \\
\mathcal R_{35} &= -9 +7x_{8} -x_{8}^{2} -8x_{8}x_{42} -4x_{8}x_{49} +24x_{42} -16x_{42}^{2} -16x_{42}x_{49} +12x_{49} \\*
&\quad -4x_{49}^{2} +x_{50} \\
\mathcal R_{36} &= -1 -4x_{50} -4x_{50}^{2} +x_{51} \\
\mathcal R_{37} &= -x_{10}^{2} -8x_{10}x_{50} -4x_{10}x_{51} -16x_{50}^{2} -16x_{50}x_{51} -4x_{51}^{2} +x_{52} \\
\mathcal R_{38} &= -x_{8} -x_{8}^{2} -4x_{8}x_{10} -4x_{8}x_{52} -4x_{10} -4x_{10}^{2} -8x_{10}x_{52} -4x_{52} \\*
&\quad -4x_{52}^{2} +x_{53} \\
\mathcal R_{39} &= -10 -12x_{53} -4x_{53}^{2} +x_{54} \\
\mathcal R_{40} &= -19 +4x_{9} -4x_{9}^{2} -8x_{9}x_{23} -16x_{9}x_{53} -8x_{9}x_{54} +4x_{23} -4x_{23}^{2} \\*
&\quad -16x_{23}x_{53} -8x_{23}x_{54} -16x_{53} -24x_{53}^{2} -16x_{53}x_{54} +6x_{54} -4x_{54}^{2} \\*
&\quad +x_{55} \\
\mathcal R_{41} &= -10 +x_{9} -25x_{9}^{2} -40x_{9}x_{23} -20x_{9}x_{55} -16x_{23}^{2} -16x_{23}x_{55} -12x_{48} \\*
&\quad -4x_{48}^{2} -4x_{55}^{2} +x_{56} \\
\mathcal R_{42} &= 21 +24x_{48} -8x_{48}^{2} -16x_{48}x_{56} -2x_{56} -4x_{56}^{2} +x_{57} \\
\mathcal R_{43} &= -5 +4x_{8} +8x_{9} +x_{10} -x_{10}^{2} -4x_{10}x_{57} -2x_{19} -4x_{57}^{2} +x_{58} +12x_{59} \\*
&\quad -4x_{59}^{2} \\
\mathcal R_{44} &= -6x_{8} -16x_{8}^{2} -64x_{8}x_{9} +8x_{8}x_{10} +16x_{8}x_{19} -12x_{9} -64x_{9}^{2} +16x_{9}x_{10} \\*
&\quad +32x_{9}x_{19} +3x_{10} -x_{10}^{2} -4x_{10}x_{19} +3x_{19} -4x_{19}^{2} +x_{59} \\
\mathcal R_{45} &= -12 +8x_{8} +16x_{9} -4x_{19} -4x_{58}^{2} +24x_{59} -8x_{59}^{2} +x_{60} \\
\mathcal R_{46} &= -1 -4x_{60} -4x_{60}^{2} +x_{61} \\
\mathcal R_{47} &= -x_{2}x_{13} +x_{62} \\
\mathcal R_{48} &= x_{12} -x_{13} -2x_{22}x_{62} -x_{63} \\
\mathcal R_{49} &= x_{12} -x_{13} -2x_{26}x_{62} -x_{64} \\
\mathcal R_{50} &= x_{12} -x_{13} -2x_{28}x_{62} -6x_{62} -x_{65} \\
\mathcal R_{51} &= x_{12} -x_{13} -2x_{30}x_{62} -x_{66} \\
\mathcal R_{52} &= x_{12} -x_{13} -2x_{32}x_{62} +2x_{62} -x_{67} \\
\mathcal R_{53} &= x_{12} -x_{13} -4x_{34}x_{62} -2x_{35}x_{62} -x_{68} \\
\mathcal R_{54} &= x_{12} -x_{13} -2x_{38}x_{62} -x_{69} \\
\mathcal R_{55} &= 2x_{28}x_{62} -2x_{41}x_{62} -6x_{62} +x_{65} -x_{70} \\
\mathcal R_{56} &= -x_{8}x_{9} +x_{8}x_{10} +x_{9} +x_{9}x_{10} -x_{10} -x_{10}^{2} +x_{71} \\
\mathcal R_{57} &= 1 +2x_{11} +x_{11}^{2} -x_{71}^{2} +x_{72} \\
\mathcal R_{58} &= -x_{12}^{2} +2x_{12}x_{13} +2x_{12}x_{74} -x_{13}^{2} -2x_{13}x_{74} -x_{72}^{2} +x_{73} \\*
&\quad -x_{74}^{2} \\
\mathcal R_{59} &= -4x_{60}x_{62} -2x_{61}x_{62} +x_{74} \\
\mathcal R_{60} &= -x_{9}x_{62} -x_{13} +x_{75}x_{76} +x_{77} \\
\mathcal R_{61} &= -1 +4x_{2} +4x_{2}x_{8} +8x_{2}x_{9} -2x_{2}x_{19} +x_{75} \\
\mathcal R_{62} &= -x_{8}x_{208} -x_{12} +x_{76} +x_{208} \\
\mathcal R_{63} &= x_{9}x_{208} -2x_{19}x_{62} +2x_{32}x_{62} +x_{67} +x_{78} \\
\mathcal R_{64} &= -x_{12}x_{63} +x_{13}x_{63} +x_{79} \\
\mathcal R_{65} &= -x_{64}x_{79} +x_{80} \\
\mathcal R_{66} &= -x_{65}x_{80} +x_{81} \\
\mathcal R_{67} &= -x_{66}x_{81} +x_{82} \\
\mathcal R_{68} &= -x_{67}x_{82} +x_{83} \\
\mathcal R_{69} &= -x_{68}x_{83} +x_{84} \\
\mathcal R_{70} &= -x_{69}x_{84} +x_{85} \\
\mathcal R_{71} &= -x_{70}x_{85} +x_{86} \\
\mathcal R_{72} &= -x_{77}x_{86} +x_{87} \\
\mathcal R_{73} &= -x_{78}x_{87} +x_{88} \\
\mathcal R_{74} &= -x_{88}^{2} +x_{89} \\
\mathcal R_{75} &= -x_{4} -2x_{4}x_{16} +2x_{16} +x_{90} \\
\mathcal R_{76} &= -x_{14}x_{90} +x_{91} \\
\mathcal R_{77} &= -2x_{1}^{2} +2x_{1}x_{12} +2x_{1}x_{13} +2x_{1}x_{91} +2x_{1}x_{92} -x_{12}^{2} -2x_{12}x_{91} \\*
&\quad -x_{13}^{2} -2x_{13}x_{92} +x_{17}x_{94} -x_{73}x_{89} -x_{91}^{2} -x_{92}^{2} \\
\mathcal R_{78} &= -x_{4}x_{15} +x_{15} -x_{15}x_{90} +x_{92} \\
\mathcal R_{79} &= -x_{95}x_{98} +x_{98} +x_{100} \\
\mathcal R_{80} &= -x_{3}x_{121} +x_{96}x_{200} -x_{99} -x_{100}x_{121} \\
\mathcal R_{81} &= x_{101}x_{200} -x_{102}x_{124} -x_{103} \\
\mathcal R_{82} &= x_{104}x_{200} -x_{105}x_{124} -x_{106} \\
\mathcal R_{83} &= x_{107}x_{200} -x_{108}x_{124} -x_{109} \\
\mathcal R_{84} &= x_{110}x_{200} -x_{111}x_{124} -x_{112} \\
\mathcal R_{85} &= x_{113}x_{201} -x_{115}x_{124} -x_{116} \\
\mathcal R_{86} &= x_{117}x_{201} -x_{118}x_{124} -x_{119} \\
\mathcal R_{87} &= -x_{17}x_{121} +x_{124} \\
\mathcal R_{88} &= x_{120} -x_{122}x_{124} -x_{123} \\
\mathcal R_{89} &= -x_{124}x_{126} +x_{125} -x_{127} \\
\mathcal R_{90} &= -x_{124}x_{129} +x_{128}x_{200} -x_{130} \\
\mathcal R_{91} &= x_{131} -2x_{131}x_{133} +x_{134} \\
\mathcal R_{92} &= -x_{131}x_{134} -x_{134} +x_{135} \\
\mathcal R_{93} &= -1 +x_{132}^{2} +x_{134}^{2} -x_{135}^{2} \\
\mathcal R_{94} &= 1 -x_{132}x_{137} -x_{136} +x_{138} \\
\mathcal R_{95} &= -1 +2x_{136} -x_{136}^{2} +2x_{136}x_{138} +x_{137}^{2} \\
\mathcal R_{96} &= 1 +x_{93}^{2} -2x_{93}x_{132} +x_{139} \\
\mathcal R_{97} &= x_{2} -x_{93}x_{182} +x_{139}x_{140} -x_{181} \\
\mathcal R_{98} &= 1 +x_{3}^{2} -2x_{3}x_{132} +x_{141} \\
\mathcal R_{99} &= -x_{3}x_{173} +x_{17} +x_{141}x_{142} -x_{172} \\
\mathcal R_{100} &= -1 +2x_{132}x_{184} -x_{143} -x_{184}^{2} \\
\mathcal R_{101} &= -1 +x_{95} +x_{143}x_{144} -x_{172} -x_{173}x_{184} \\
\mathcal R_{102} &= -x_{132}x_{146} -x_{145} +x_{147} \\
\mathcal R_{103} &= -1 -x_{145}^{2} +2x_{145}x_{147} +x_{146}^{2} \\
\mathcal R_{104} &= 1 -x_{132}x_{149} -x_{148} +x_{150} \\
\mathcal R_{105} &= -1 +2x_{148} -x_{148}^{2} +2x_{148}x_{150} +x_{149}^{2} \\
\mathcal R_{106} &= -x_{132}x_{152} -x_{151} +x_{153} \\
\mathcal R_{107} &= -1 -x_{151}^{2} +2x_{151}x_{153} +x_{152}^{2} \\
\mathcal R_{108} &= -x_{132}x_{155} -x_{154} +x_{156} \\
\mathcal R_{109} &= -1 -x_{154}^{2} +2x_{154}x_{156} +x_{155}^{2} \\
\mathcal R_{110} &= -x_{132}x_{158} -x_{157} +x_{159} \\
\mathcal R_{111} &= -1 -x_{157}^{2} +2x_{157}x_{159} +x_{158}^{2} \\
\mathcal R_{112} &= -x_{132}x_{161} -x_{160} +x_{162} \\
\mathcal R_{113} &= -1 -x_{160}^{2} +2x_{160}x_{162} +x_{161}^{2} \\
\mathcal R_{114} &= 1 -x_{132}x_{164} -x_{163} +x_{165} \\
\mathcal R_{115} &= -1 +2x_{163} -x_{163}^{2} +2x_{163}x_{165} +x_{164}^{2} \\
\mathcal R_{116} &= -x_{132}x_{167} -x_{166} +x_{168} \\
\mathcal R_{117} &= -1 -x_{166}^{2} +2x_{166}x_{168} +x_{167}^{2} \\
\mathcal R_{118} &= -x_{132}x_{170} -x_{169} +x_{171} \\
\mathcal R_{119} &= -1 -x_{169}^{2} +2x_{169}x_{171} +x_{170}^{2} \\
\mathcal R_{120} &= 1 -x_{132}x_{173} -x_{172} +x_{174} \\
\mathcal R_{121} &= -1 +2x_{172} -x_{172}^{2} +2x_{172}x_{174} +x_{173}^{2} \\
\mathcal R_{122} &= -x_{132}x_{176} -x_{175} +x_{177} \\
\mathcal R_{123} &= -1 -x_{175}^{2} +2x_{175}x_{177} +x_{176}^{2} \\
\mathcal R_{124} &= -x_{132}x_{179} -x_{178} +x_{180} \\
\mathcal R_{125} &= -1 -x_{178}^{2} +2x_{178}x_{180} +x_{179}^{2} \\
\mathcal R_{126} &= 1 -x_{132}x_{182} -x_{181} +x_{183} \\
\mathcal R_{127} &= -1 +2x_{181} -x_{181}^{2} +2x_{181}x_{183} +x_{182}^{2} \\
\mathcal R_{128} &= -5 -2x_{0} -3x_{0}^{2} -2x_{0}x_{2} -x_{1}^{2} -2x_{2}^{2} -x_{3}^{2} -x_{4}^{2} +2x_{8} \\*
&\quad -x_{8}^{2} -x_{9}^{2} -x_{10}^{2} -x_{11}^{2} -x_{12}^{2} -x_{13}^{2} -x_{14}^{2} \\*
&\quad -x_{15}^{2} -4x_{16}^{2} -x_{17}^{2} -2x_{93}^{2} +2x_{95} -2x_{95}^{2} -x_{96}^{2} \\*
&\quad -x_{97}^{2} -x_{101}^{2} -x_{104}^{2} -x_{107}^{2} -x_{110}^{2} -x_{113}^{2} \\*
&\quad -x_{114}^{2} -x_{117}^{2} -x_{120}^{2} -x_{121}^{2} -x_{125}^{2} -x_{128}^{2} \\*
&\quad +x_{131} -x_{184}^{2} \\
\mathcal R_{129} &= 2x_{3} -4x_{3}x_{8} -4x_{3}x_{9} -4x_{3}x_{10} -4x_{3}x_{11} -4x_{3}x_{12} -4x_{3}x_{13} -4x_{3}x_{14} \\*
&\quad -4x_{3}x_{15} -8x_{3}x_{16} +x_{184} \\
\mathcal R_{130} &= -464 +x_{0} -x_{185}^{2} \\
\mathcal R_{131} &= -x_{132}x_{187} +x_{196} \\
\mathcal R_{132} &= -1 +x_{186}^{2} +x_{187}^{2} -x_{196}^{2} \\
\mathcal R_{133} &= -x_{189}x_{190} +x_{197} \\
\mathcal R_{134} &= -1 +x_{188}^{2} +x_{189}^{2} -x_{197}^{2} \\
\mathcal R_{135} &= -x_{179}^{2} +x_{198} \\
\mathcal R_{136} &= x_{187} -2x_{192}x_{198} -x_{198} \\
\mathcal R_{137} &= -1 -4x_{187} -8x_{187}x_{191} +x_{190} \\
\mathcal R_{138} &= -x_{132} -x_{186}x_{193} +x_{190} \\
\mathcal R_{139} &= -x_{180} -x_{186}x_{194} +x_{188} \\
\mathcal R_{140} &= -2x_{16} -4x_{179}x_{195} +x_{189} \\
\mathcal R_{141} &= 9 +2x_{0}^{2} +2x_{0}x_{2} -2x_{0}x_{137} -2x_{0}x_{146} +x_{1}^{2} -2x_{1}x_{173} +x_{2}^{2} \\*
&\quad -2x_{2}x_{146} +x_{4}^{2} -2x_{4}x_{164} -2x_{8} +x_{8}^{2} -2x_{8}x_{149} +x_{9}^{2} -2x_{9}x_{152} \\*
&\quad +x_{10}^{2} -2x_{10}x_{155} +x_{11}^{2} -2x_{11}x_{158} +x_{12}^{2} -2x_{12}x_{161} \\*
&\quad +x_{13}^{2} -2x_{13}x_{167} +x_{14}^{2} -2x_{14}x_{170} +x_{15}^{2} -2x_{15}x_{176} \\*
&\quad -4x_{16}x_{179} +x_{93}^{2} -2x_{93}x_{182} -x_{132}x_{199} +x_{138}^{2} -2x_{147} \\*
&\quad +x_{147}^{2} +2x_{149} +x_{150}^{2} -2x_{153} +x_{153}^{2} -2x_{156} +x_{156}^{2} \\*
&\quad -2x_{159} +x_{159}^{2} -2x_{162} +x_{162}^{2} +x_{165}^{2} -2x_{168} +x_{168}^{2} \\*
&\quad -2x_{171} +x_{171}^{2} +x_{174}^{2} -2x_{177} +x_{177}^{2} +2x_{179}^{2} -2x_{180} \\*
&\quad +x_{180}^{2} +x_{183}^{2} +x_{199} -x_{202}^{2} -x_{203}^{2} -x_{204} \\*
&\quad -x_{204}^{2} \\
\mathcal R_{142} &= -x_{97}x_{200} +x_{121} \\
\mathcal R_{143} &= -x_{114}x_{201} +x_{121} \\
\mathcal R_{144} &= -1 -4x_{16}^{2} -x_{137}^{2} -x_{146}^{2} -x_{149}^{2} -x_{152}^{2} -x_{155}^{2} \\*
&\quad -x_{158}^{2} -x_{161}^{2} -x_{164}^{2} -x_{167}^{2} -x_{170}^{2} -x_{173}^{2} \\*
&\quad -x_{176}^{2} +x_{179}^{2} -x_{182}^{2} -x_{202}^{2} -x_{203}^{2} -x_{204} \\*
&\quad -x_{204}^{2} \\
\mathcal R_{145} &= -1 -4x_{99}^{2} -4x_{103}^{2} -4x_{106}^{2} -4x_{109}^{2} -4x_{112}^{2} -4x_{116}^{2} \\*
&\quad -4x_{119}^{2} +x_{121}^{2} -4x_{123}^{2} -4x_{127}^{2} -4x_{130}^{2} -x_{205}^{2} \\*
&\quad -x_{206}^{2} -x_{207} -x_{207}^{2} \\
\mathcal R_{146} &= -x_{2}x_{12} +x_{208} \\
\mathcal R_{147} &= x_{95}x_{210} -2x_{96}x_{146} -2x_{101}x_{149} -2x_{104}x_{152} -2x_{107}x_{155} -2x_{110}x_{158} \\*
&\quad -2x_{113}x_{161} -2x_{117}x_{167} -2x_{120}x_{170} -2x_{125}x_{176} -2x_{128}x_{179} -2x_{132}x_{210} \\*
&\quad +2x_{146}x_{147} +2x_{149} +2x_{149}x_{150} +2x_{152}x_{153} +2x_{155}x_{156} +2x_{158}x_{159} \\*
&\quad +2x_{161}x_{162} +2x_{167}x_{168} +2x_{170}x_{171} +2x_{176}x_{177} +2x_{179}x_{180} +x_{209} \\
\mathcal R_{148} &= -10 -x_{95}x_{209} +x_{96}^{2} -2x_{96}x_{145} +x_{101}^{2} -2x_{101}x_{148} +x_{104}^{2} \\*
&\quad -2x_{104}x_{151} +x_{107}^{2} -2x_{107}x_{154} +x_{110}^{2} -2x_{110}x_{157} +x_{113}^{2} \\*
&\quad -2x_{113}x_{160} +x_{117}^{2} -2x_{117}x_{166} +x_{120}^{2} -2x_{120}x_{169} +x_{125}^{2} \\*
&\quad -2x_{125}x_{175} +x_{128}^{2} -2x_{128}x_{178} +2x_{145}x_{147} +2x_{148} +2x_{148}x_{150} \\*
&\quad +2x_{151}x_{153} +2x_{154}x_{156} +2x_{157}x_{159} +2x_{160}x_{162} +2x_{166}x_{168} +2x_{169}x_{171} \\*
&\quad +2x_{175}x_{177} +2x_{178}x_{180} +x_{210} \\
\mathcal R_{149} &= x_{95}x_{212} -2x_{97}x_{137} -2x_{114}x_{164} -2x_{121}x_{173} -2x_{132}x_{212} +2x_{137} \\*
&\quad +2x_{137}x_{138} +2x_{164} +2x_{164}x_{165} +2x_{173} +2x_{173}x_{174} +x_{211} -x_{212} \\
\mathcal R_{150} &= -2 -x_{95}x_{211} +x_{97}^{2} -2x_{97}x_{136} +x_{114}^{2} -2x_{114}x_{163} +x_{121}^{2} \\*
&\quad -2x_{121}x_{172} +2x_{136} +2x_{136}x_{138} +2x_{163} +2x_{163}x_{165} +x_{172}^{2} \\*
&\quad -x_{173}^{2} +x_{211} +x_{212} \\
\mathcal R_{151} &= -1 -x_{2} -2x_{2}x_{213} +x_{4} \\
\mathcal R_{152} &= -464x_{0}^{2} +x_{93} \\
\end{align*}
\endgroup


\end{document}
